\let\negmedspace\undefined
\let\negthickspace\undefined
\documentclass[journal]{IEEEtran}
\usepackage[a5paper, margin=10mm, onecolumn]{geometry}
%\usepackage{lmodern} % Ensure lmodern is loaded for pdflatex
\usepackage{tfrupee} % Include tfrupee package

\setlength{\headheight}{1cm} % Set the height of the header box
\setlength{\headsep}{0mm}     % Set the distance between the header box and the top of the text

\usepackage{gvv-book}
\usepackage{gvv}
\usepackage{cite}
\usepackage{amsmath,amssymb,amsfonts,amsthm}
\usepackage{algorithmic}
\usepackage{graphicx}
\usepackage{textcomp}
\usepackage{xcolor}
\usepackage{txfonts}
\usepackage{listings}
\usepackage{enumitem}
\usepackage{mathtools}
\usepackage{gensymb}
\usepackage{comment}
\usepackage[breaklinks=true]{hyperref}
\usepackage{tkz-euclide} 
\usepackage{listings}
% \usepackage{gvv}                                        
\def\inputGnumericTable{}                                 
\usepackage[latin1]{inputenc}                                
\usepackage{color}                                            
\usepackage{array}                                            
\usepackage{longtable}                                       
\usepackage{calc}                                             
\usepackage{multirow}                                         
\usepackage{hhline}                                           
\usepackage{ifthen}                                           
\usepackage{lscape}
\begin{document}

\bibliographystyle{IEEEtran}
\vspace{3cm}

\title{10.5.1}
\author{EE25BTECH11033 - Kavin}
% \maketitle
% \newpage
% \bigskip
{\let\newpage\relax\maketitle}

\renewcommand{\thefigure}{\theenumi}
\renewcommand{\thetable}{\theenumi}
\setlength{\intextsep}{10pt} % Space between text and floats
\textbf{Question}:\\
Let $ABC$ be a right triangle in which $AB = 6 cm, BC = 8 cm$ and $\angle B = 90\degree$. $BD$ is the perpendicular from $\vec{B}$ on $AC$. The circle through $\vec{B}, \vec{C}, \vec{D}$ is drawn. Construct the tangents from $\vec{A}$ to this circle.\\    
\bigskip

\textbf{Solution}:\\

Let,
\begin{align}
    \vec{B} = \myvec{0\\0} , \vec{A}=\myvec{0\\6} \text{ and } \vec{C} = \myvec{8\\0}
\end{align}

Let the equation of the line passing through $\vec{A}$ and $\vec{C}$ be,
\begin{align}
    \vec{n}^\top\vec{x}=c
\end{align}
The direction vector of line $AC$ is given by,
\begin{align}
\vec{m}=\vec{A}-\vec{C} \equiv \myvec{-4\\3}
\end{align}
Therefore, the normal vector of the desired line is given by,
\begin{align}
    \vec{n}=\myvec{3\\4}
\end{align}
\begin{align}
    c = \vec{n}^\top\vec{A} = \myvec{3&4}\myvec{0\\6} = 24
\end{align}
Let $\vec{D}$ be the foot of the perpendicular from $\vec{B}$ to the line
\begin{align}
    \myvec{\vec{m} & \vec{n}}^\top\vec{D} &= 
	   \myvec{\vec{m}^\top\vec{B}\\c}
\end{align}
\begin{align}
    \myvec{-4 & 3\\3 & 4}\vec{D}=\myvec{\vec{m}^\top\vec{B}\\24}
\end{align}
\begin{align}
    \vec{m}^\top\vec{B} =  \myvec{-4 & 3}\myvec{0 \\ 0} = 0
\end{align}
\begin{align}
    \implies\myvec{-4 & 3\\3 & 4}\vec{D}=\myvec{0\\24}
\end{align}
The augmented matrix  is given by,
\begin{align}
    \augvec{2}{1}{-4 & 3 & 0\\3 & 3 & 24}
\end{align}
\begin{align}
    R_1\rightarrow \frac{-1}{4}R_1\implies\augvec{2}{1}{1 & -3/4 & 0\\3 & 3 & 24}
\end{align}
\begin{align}
    R_2 \rightarrow R_2-3R_1 \implies\augvec{2}{1}{1 & -3/4 & 0\\0 & 25/4 & 24}
\end{align}
\begin{align}
    R_2 \rightarrow \frac{4}{25}R_2\implies\augvec{2}{1}{1 & -3/4 & 0\\0 & 1 & 96/25}
\end{align}
\begin{align}
    R_1 \rightarrow R_1 + \frac{3}{4}R_2\implies\augvec{2}{1}{1 & 0 & 72/25\\0 & 1 & 96/25}
\end{align}
\begin{align}
    \vec{D} = \myvec{72/25 \\ 96/25}
\end{align}
To find the circle through $\vec{B}, \vec{C}, \vec{D}$ the parameters are given by
\begin{align}
\myvec{
 2 \vec{B} & 2 \vec{C} & 2\vec{D}
 \\
 1 & 1 & 1
 }^\top 
	\myvec{\vec{u} \\ f}
	=
-\myvec{\norm{\vec{B}}^2 \\\norm{\vec{C}}^2 \\\norm{\vec{D}}^2 	}
\end{align}
\begin{align}
    \myvec{0 & 0 & 1\\16 & 0 & 1\\144/25 & 192/25 & 1}\myvec{\vec{u} \\ f} = \myvec{0\\-64\\-576/25}
\end{align}
The augmented matrix  is given by,
\begin{align}
    \augvec{3}{1}{0 & 0 & 1 & 0\\16 & 0 & 1 & -64\\144/25 & 192/25 & 1 & -576/25}
\end{align}
\begin{align}
    R_1\leftrightarrow R_2\implies\augvec{3}{1}{16 & 0 & 1 & -64\\0 & 0 & 1 & 0\\144/25 & 192/25 & 1 & -576/25}
\end{align}
\begin{align}
    R_1 \rightarrow \frac{1}{16}R_1 \implies\augvec{3}{1}{1 & 0 & 1/16 & -4\\0 & 0 & 1 & 0\\144/25 & 192/25 & 1 & -576/25}
\end{align}
\begin{align}
    R_3 \rightarrow R_3 - \frac{144}{25}R_1\implies\augvec{3}{1}{1 & 0 & 1/16 & -4\\0 & 0 & 1 & 0\\0 & 192/25 & 16/25 & -504/25}
\end{align}
\begin{align}
    R_2\leftrightarrow R_3\implies\augvec{3}{1}{1 & 0 & 1/16 & -4\\0 & 192/25 & 16/25 & -504/25\\0 & 0 & 1 & 0}
\end{align}
\begin{align}
    R_2 \rightarrow \frac{25}{192}R_2 \implies\augvec{3}{1}{1 & 0 & 1/16 & -4\\0 & 1 & 1/12 & 0\\0 & 0 & 1 & 0}
\end{align}
\begin{align}
    R_2 \rightarrow R_2 - \frac{1}{12}R_3\ ,\ R_1 \rightarrow R_1 - \frac{1}{16}R_3 \implies\augvec{3}{1}{1 & 0 & 0 & -4\\0 & 1 & 0 & 0\\0 & 0 & 1 & 0}
\end{align}

\begin{align}
    \implies\vec{u} = \myvec{-4 \\ 0} \text{ and } f=0
\end{align}\\

Let $\vec{n}^\top \vec{x}=c$ be the equation of tangent to the circle.\\
Also,
\begin{align}
    \vec{n}^\top \vec{A}=c
\end{align}
and,
\begin{align}
    \frac{\abs{\vec{n}^\top(-\vec{u})-c}}{\norm{\vec{n}}} = r
\end{align}
where
\begin{align*}
    r=\sqrt{\norm{\vec{u}}^2-f}
\end{align*}
\begin{align}
    \implies \vec{n}^\top \vec{u} + c= \pm \norm{\vec{n}}\sqrt{\norm{\vec{u}}^2-f}
\end{align}
\begin{align}
    \implies  \vec{n}^\top=\brak{\vec{u} + \vec{A}} \pm \norm{\vec{n}}\sqrt{\norm{\vec{u}}^2-f}
\end{align}
substitute$ \vec{n} = \myvec{-m\\1}$ where m is the slope of the line.
\begin{align}
    \implies  \myvec{-m & 1}\myvec{-4\\6} = \pm4\sqrt{1+m^2}
\end{align}
\begin{align}
    m=-5/12 \implies \vec{n}=\myvec{5\\12}
\end{align}
\begin{align}
    c=\vec{n}^\top \vec{A} =72
\end{align}\\
Therefore, the line $\myvec{5 & 12}\vec{x} = 72$ and the y-axis are the two tangents.\\
\bigskip 

\begin{figure}[H]
\begin{center}
\includegraphics[width=0.7\columnwidth]{figs/fig.png}
\end{center}
\label{fig:Fig1}
\end{figure}
\end{document}


